% Pista ana-24i-q1 · Anàlisi
% Procedència: PAU Catalunya 2024 · Convocatòria de juny (incidències) · Sèrie 5 · Qüestió 1
\documentclass[11pt,a4paper]{article}
\usepackage{pista}
\pistainfo{Catalunya 2024}{Convocatòria de juny (incidències)}{Sèrie 5}

\begin{document}

\section*{\textcolor{blauPista}{Qüestió 1}}

\noindent Considereu la funció $f(x) = -2 + 10\,(x-1)\ln x$, definida per a $x > 0$.

\subsection*{\textit{a)} \normalfont Comproveu que $f(x)$ té una arrel a l'interval $[1, 1{,}5]$ i busqueu un interval d'una dècima de longitud que també contingui aquesta mateixa arrel. \hspace*{\fill}\textnormal{[0,75~punts]}}

\pista{Aplica el teorema de Bolzano: si $f$ és contínua en l'interval i canvia de signe entre els extrems, hi ha alguna arrel.
\begin{itemize}
\item Calcula $f(1)$ i $f(1{,}5)$ i comprova-ho.
\item Per trobar l'interval d'amplitud $0{,}1$, prova valors intermedis ($1{,}1$, $1{,}2$, etc.) fins a localitzar un canvi de signe.
\end{itemize}}

\subsection*{\textit{b)} \normalfont Sense calcular els punts crítics, justifiqueu que $f(x)$ és decreixent a l'interval $(0, 1)$ i creixent a $(1, +\infty)$. Quins màxims i mínims té aquesta funció? \hspace*{\fill}\textnormal{[1~punt]}}

\pista{«Sense calcular els punts crítics» vol dir que no has de resoldre $f'(x) = 0$, però sí que pots derivar.
\begin{itemize}
\item Calcula $f'(x)$ i estudia'n el signe a $(0, 1)$ i a $(1, +\infty)$: quin signe tenen $\ln x$ i $x - 1$ a cada interval?
\item Del creixement i el decreixement, dedueix on hi ha el mínim.
\item Decideix també si hi ha algun màxim o no n'hi ha.
\end{itemize}}

\subsection*{\textit{c)} \normalfont Calculeu $\displaystyle \lim_{x \to 0^{+}} f(x)$ i $\displaystyle \lim_{x \to +\infty} f(x)$, i feu un esbós de la gràfica d'aquesta funció. \hspace*{\fill}\textnormal{[0,75~punts]}}

\pista{Segueix els passos de l'estudi d'una funció. Alguns ja els tens dels apartats anteriors.
\begin{enumerate}
\item \textit{Domini i asímptotes verticals.} El domini és $x > 0$. Quan $x \to 0^{+}$, $(x - 1) \to -1$ i $\ln x \to -\infty$: dedueix-ne el límit directament, sense eines avançades. Si és infinit, $x = 0$ és una asímptota vertical.
\item \textit{Límits a l'infinit.} Pel domini, només cal $x \to +\infty$: els dos factors de $(x - 1)\ln x$ tendeixen a $+\infty$. (La constant $-2$ no influeix en cap dels dos límits.)
\item \textit{Talls amb els eixos.} Amb l'eix d'ordenades no n'hi ha, perquè $x = 0$ no és del domini. Amb l'eix d'abscisses, tens una arrel localitzada a l'apartat \textit{a)}. Fixa't en el signe de $f(1)$: quantes vegades talla la gràfica l'eix d'abscisses?
\item \textit{Creixement, decreixement i extrems relatius.} Ho tens de l'apartat \textit{b)}.
\item \textit{Concavitat i punts d'inflexió.} Calcula $f''(x)$: quin signe té per a $x > 0$?
\item \textit{Taula de valors} (opcional): uns quants punts t'ajudaran a situar la corba.
\item \textit{Esbós de la gràfica.} Combina-ho tot.
\end{enumerate}}

\end{document}
